\documentclass[11pt,a4paper]{article}
\usepackage{operationsform}
\begin{document}
\OPStitle{複数媒体間の内容整合チェックリスト}{TPL-CHK-011}{0.2.0}
\OPSsection{1}{確認項目}
\OPScheck{1}{長文、ニュース、FAQ、SNS、動画、音声で中心的な主張が一致する}
\OPScheck{2}{対象となる人数、期間、適用条件の記載に重要な矛盾がない}
\OPScheck{3}{直接引用、要約、翻訳に同じ原記録IDが付いている}
\OPScheck{4}{画像・キャプション・字幕が本文より強く断定していない}
\OPScheck{5}{短縮版で除かれた限定条件へ読者が参照できる経路がある}
\OPScheck{6}{媒体ごとの出典と権利、解禁条件を確認した}
\OPScheck{7}{ある媒体の改訂が他媒体へ反映されたか確認した}
\OPScheck{8}{版・確認担当・未決項目を媒体別に残した}
\OPSsection{2}{点検結果の記録}
\OPSblankshort{対象案件・対象版}
\OPSblankshort{担当者と点検日}
\OPSblanklong{未解決事項と対応}
\par\Needspace{4\baselineskip}\noindent\textbf{確認に用いた資料ID}\par\vspace{1mm}
{\fontsize{8.5}{13}\selectfont
\begin{tabularx}{\linewidth}{@{}XXXX@{}}
\textbf{対象・項目} & \textbf{内容・参照} & \textbf{担当・状態} & \textbf{備考}\\[2mm]
\strut & \strut & \strut & \strut\\[5mm]
\strut & \strut & \strut & \strut\\[5mm]
\end{tabularx}
}\vspace{2mm}\par
\end{document}
